\textbf{Время:} 01:30

\textbf{Тезис:} STAR регистрирует пионы из столкновений золота с золотом, чтобы исследовать область их испускания после сжатия и расширения вещества.

\textbf{На экране:} слева направо по схеме — пучок Au, неподвижная мишень Au, выходящие пионы и TPC + TOF; затем две фразы справа о физической задаче; в конце нижняя подпись с определением энергии √s\_NN.

\textbf{Текст:} Мы разобрали, как измерить импульс и выбрать пионы. Теперь свяжем эти приборы с конкретным экспериментом. На схеме пучок ядер золота попадает в неподвижную золотую мишень. Возникающие заряженные частицы регистрирует STAR; TPC и TOF нам уже знакомы.

Справа сформулирована задача: исследовать сжатие и последующее расширение ядерного вещества. После столкновения вещество расширяется и становится разреженным. Зарегистрированные пионы несут информацию об области, из которой они окончательно вышли. Её характерные размеры будем определять по корреляциям.

В названии указан режим при трёх ГэВ. Нижняя подпись поясняет, что это энергия пары сталкивающихся нуклонов в их системе центра масс. Это не энергия всего ядра и не энергия пучка в лаборатории. Для дальнейшего рассказа важно, что все показанные затем данные STAR относятся к одному определённому режиму столкновений.

Следующий шаг уже статистический: отдельные пионы измерены, теперь из них нужно составить пары и построить корреляционную функцию.

\textbf{Переход:} Сравним пары из одного столкновения с фоном, собранным из разных похожих столкновений.

\textbf{Справка:} Au — золото. √s\_NN — полная энергия пары нуклонов в их системе центра масс; в данном исследовании 3 ГэВ. Схема продольная, без масштаба; положение мишени и геометрия регистрации показаны качественно. При этой энергии существенна остановка исходных нуклонов и интерес представляет плотное ядерное вещество. Фазовую диаграмму и химический потенциал в основной рассказ не вводить. Образование кварк-глюонной плазмы при 3 ГэВ по этим пионным корреляциям не установлено; термодинамический контекст оставлен в резерве. А. Краева (STAR), Universe 10, 188 (2024), введение и §2.2; DOI: \url{https://doi.org/10.3390/universe10040188}. Дополнительная схема установки находится в резерве 38, фазовая диаграмма — в резерве 40.
